\uuid{XD5y}
\titre{Système linéaire à un paramètre/Méthode du Pivot de Gauss}
\auteur{hanani ; mzari}
\organisation{exo7}
\licence{CC BY-NC-SA 4.0}
\datecreate{2020-01-20}
\niveau{L1}
\module{Algèbre}
\chapitre{Déterminant, système linéaire}
\sousChapitre{Système linéaire, rang}
\difficulte{3}
\theme{Systèmes d'équations linéaires, Système linéaire à un paramètre, Méthode du Pivot de Gauss}
\exo7qcmid{lille-2.systemes.011}

\contenu{
\begin{question}
\qclassification{200.03}
\qtags{motcle=Système linéaire à un paramètre/Méthode du Pivot de Gauss}
On considère le système d'équations, d'inconnue $(x,y,z)\in \Rr^3$ et de paramètre un réel $m$ :
$$(\mathtt{S}) \; \left\{\begin{array}{rcc}
x+y+z&=&-1\\
x+2y+3z&=&1\\
2x+3y+4z&=&m.\end{array}\right.$$
Quelles sont les assertions vraies ?
\begin{answers}
    \bad{$(\mathtt{S}) \Leftrightarrow
\left\{\begin{array}{rcc}
x+y+z&=&-1\\
y+2z&=&m.\\
\end{array}\right.$}
    \bad{Pour tout réel $m$, $(\mathtt{S})$ admet une solution.}
    \good{Si $m=1$, $(\mathtt{S})$ n'admet pas de solution.}
    \good{Si $m=0$, l'ensemble des solutions de $(\mathtt{S})$ est une droite.}
\end{answers}
\begin{explanations}
$$(\mathtt{S}) \Leftrightarrow
\left\{\begin{array}{rcc}
x+y+z&=&-1\\
y+2z&=&2\\
y+2z&=&m+2\\
\end{array}\right. \Leftrightarrow
\left\{\begin{array}{rcc}
x+y+z&=&-1\\
y+2z&=&2\\
0&=&m.\\
\end{array}\right.$$
On en déduit que si $m\neq 0$, $(\mathtt{S})$ n'admet pas de solution et que si $m=0$, l'ensemble des solutions de
$(\mathtt{S})$ est la droite :
$\{(-3+z,2-2z,z); \; z\in \Rr\}$.
\end{explanations}
\end{question}
}
